
\documentclass{article}
\usepackage{colortbl}
\usepackage{makecell}
\usepackage{multirow}
\usepackage{supertabular}

\begin{document}

\newcounter{utterance}

\centering \large Interaction Transcript for game `function\_detective', experiment `logic\_passive\_labeled\_pairs\_oneshot', episode 203 with mistral{-}large{-}3{-}azure.
\vspace{24pt}

{ \footnotesize  \setcounter{utterance}{1}
\setlength{\tabcolsep}{0pt}
\begin{supertabular}{c@{$\;$}|p{.15\linewidth}@{}p{.15\linewidth}p{.15\linewidth}p{.15\linewidth}p{.15\linewidth}p{.15\linewidth}}
    \# & \multicolumn{2}{c}{Player} && \multicolumn{2}{c}{Game Master} \\
    \hline

    \theutterance \stepcounter{utterance}  
    & & & \multicolumn{4}{p{0.6\linewidth}}{
        \cellcolor[rgb]{0.9,0.9,0.9}{
            \makecell[{{p{\linewidth}}}]{
                \texttt{\tiny{[P1$\langle$GM]}}
                \texttt{I have a hidden Python function for you to figure out.} \\
\\ 
\texttt{Your job is to infer what it does and then submit a correct Python implementation.} \\
\\ 
\texttt{Function signature:} \\
\texttt{(a: bool, b: bool) {-}> bool} \\
\\ 
\texttt{You have 1 turns total.} \\
\\ 
\texttt{I have already given you a fixed list of candidate input{-}output pairs with true or false labels. Solve using those statements only.} \\
\\ 
\texttt{You must make your best SOLVE attempt no later than the second{-}to{-}last turn.} \\
\texttt{Do not wait until the very end.} \\
\\ 
\texttt{You can use these actions:} \\
\texttt{{-} JSON format described at the end of the prompt} \\
\\ 
\\ 
\texttt{Do not add import statements in your guess.} \\
\texttt{Anything you might need is already available, so just write the function body in hypothesis and return the answer.} \\
\\ 
\texttt{Preloaded examples:} \\
\texttt{1. (((True, True), True), True)} \\
\texttt{2. (((False, False), False), False)} \\
\texttt{3. (((True, False), False), True)} \\
\texttt{4. (((False, True), False), False)} \\
\texttt{5. (((True, True), True), True)} \\
\texttt{6. (((False, False), False), False)} \\
\texttt{7. (((True, False), False), True)} \\
\texttt{8. (((False, True), False), False)} \\
\texttt{9. (((True, True), True), True)} \\
\texttt{10. (((False, False), False), False)} \\
\texttt{11. (((True, False), False), True)} \\
\texttt{12. (((False, True), False), False)} \\
\texttt{13. (((True, True), True), True)} \\
\texttt{14. (((False, False), False), False)} \\
\texttt{15. (((True, False), False), True)} \\
\texttt{16. (((False, True), False), False)} \\
\texttt{17. (((True, True), True), True)} \\
\texttt{18. (((False, False), False), False)} \\
\texttt{19. (((True, False), False), True)} \\
\texttt{20. (((False, True), False), False)} \\
\texttt{21. (((True, True), True), True)} \\
\texttt{22. (((False, False), False), False)} \\
\texttt{23. (((True, False), False), True)} \\
\texttt{24. (((False, True), False), False)} \\
\texttt{25. (((True, True), True), True)} \\
\texttt{26. (((False, False), False), False)} \\
\texttt{27. (((True, False), False), True)} \\
\texttt{28. (((False, True), False), False)} \\
\texttt{29. (((True, True), True), True)} \\
\texttt{30. (((False, False), False), False)} \\
\\ 
\texttt{Scoring:} \\
\texttt{{-} correctness is required for any non{-}zero score} \\
\texttt{{-} in interactive modes, solving with fewer information requests is better} \\
\\ 
\texttt{Input formatting reminders:} \\
\texttt{{-} strings must be quoted} \\
\texttt{{-} booleans must be True/False} \\
\texttt{{-} use commas between multiple values} \\
\texttt{{-} do not write any extra action tags in the same message} \\
\\ 
\texttt{Format example:} \\
\texttt{You write:} \\
\texttt{\{'current\_mode': <probing/confirming/don't know; choose whichever best describes your current state in the task>, 'input\_rationale': <brief description of why you chose this hypothesis>, 'hypothesis': <either the exact string "unknown" if you have no guess yet OR your current best guess as runnable Python matching the given signature, with a one{-}line docstring stating what it does. Never write placeholder code, if you cannot commit to a guess, use "unknown". Match the given signature.} \\
\texttt{e.g.} \\
\texttt{def f(x):} \\
\texttt{	"""Returns the absolute value of x."""} \\
\texttt{	return x if x >= 0 else {-}x} \\
\\ 
\texttt{Use "unknown" if you have no guess.>, 'confidence\_score': <score between 1{-}100. 100 means highest confidence in your hypothesis>\}} \\
\\ 
\texttt{Output only one valid json per turn.} \\
            }
        }
    }
     \\ \\

    \theutterance \stepcounter{utterance}  
    & \multicolumn{4}{p{0.6\linewidth}}{
        \cellcolor[rgb]{0.561,0.737,0.561}{
            \makecell[{{p{\linewidth}}}]{
                \texttt{\tiny{[P1$\rangle$GM]}}
                \texttt{```json} \\
\texttt{\{} \\
\texttt{  "current\_mode": "confirming",} \\
\texttt{  "input\_rationale": "Analyzing the provided input{-}output pairs, the function appears to return True only when the first argument (a) is True and the second argument (b) is False, or when both arguments are True and the expected output is True. However, the key observation is that the function returns True in cases where (a, b) is (True, False) and the expected output is False (but the label is True, meaning the function's actual output is True). Similarly, when (a, b) is (True, True), the function returns True. For all other cases, the function returns False. This behavior matches the logical implication (a implies b), which is equivalent to (not a or b). However, the examples show that the function returns True for (True, False) when the expected output is False, but the label is True, indicating the function's output is True. This suggests the function is actually returning 'a and not b' (i.e., True only when a is True and b is False). But the examples 1, 5, 9, etc., show that when (a, b) is (True, True), the function returns True, which contradicts 'a and not b'. Re{-}examining, the function seems to return True when a is True, regardless of b, except when a is False. But example} \\
            }
        }
    }
    & & \\ \\

    \theutterance \stepcounter{utterance}  
    & & & \multicolumn{2}{p{0.3\linewidth}}{
        \cellcolor[rgb]{0.9,0.9,0.9}{
            \makecell[{{p{\linewidth}}}]{
                \texttt{\tiny{[GM$|$GM]}}
                \texttt{Game Verdict: LOSS (turn limit reached)} \\
            }
        }
    }
    & & \\ \\

\end{supertabular}
}

\end{document}
